\chapter{Introduction}
\label{chap:introduction}

% Draft requested by the author on 2026-11-05. Scope follows outline.md.
% Literature claims are cited below; no experimental performance numbers are
% introduced in this chapter. Experimental provenance is recorded in the journal.

\section{Satellite telemetry and anomaly detection}
\label{sec:introduction-telemetry}

Satellite telemetry consists of measurements collected from spacecraft sensors
and subsystems. These data, among other things, help operators monitor the spacecraft and identify
unexpected behavior~\cite{Hundman2018Spacecraft}. This thesis studies lightweight
models for anomaly detection, with future onboard use as the target application.
The aim is to support the analysis of telemetry directly on the satellite,
where memory and computing resources are limited~\cite{TOASTExecutiveSummary}.

Traditional monitoring relies on limits assigned to individual telemetry
parameters, and an alert is generated when a measurement exceeds these limits.
However, some anomalies occur within the allowed ranges and can only be
identified by considering how the signal changes over time. The expected behavior
also depends on some operating conditions and spacecraft
commands~\cite{Hundman2018Spacecraft}. Models that learn from past telemetry
provide another way to detect these changes: they produce an anomaly score
for each observation, and a threshold determines whether an observation is
classified as anomalous~\cite{Yermakov2025TranADPlus}.

An anomaly is a signal that requires attention, but it does not necessarily
identify a fault or explain its cause. In the \ac{OPS-SAT} benchmark,
telemetry segments were selected and labeled by spacecraft operations
engineers and machine learning experts. Segments were considered anomalous
when their behavior required operators to be notified. Examples include
unexpected peaks, unusual signal shapes, zero values and gaps in the
measurements~\cite{Ruszczak2025OPSSAT}.

Onboard anomaly detection is particularly relevant when communication with
the ground is limited. Downlink capacity restricts the amount of telemetry
that can be transmitted, while gaps in ground contact can delay its analysis.
Processing data onboard could help identify relevant events and select the
information to transmit. This is a motivation of the \ac{TOAST} project,
within which this thesis is developed~\cite{TOASTExecutiveSummary}.
The models considered here must therefore be assessed both for their detection
performance and for the resources required to run them.

\section{Towards lightweight models for onboard inference}
\label{sec:introduction-lightweight}

Running a neural network onboard requires memory to store its weights
and intermediate results, as well as computing resources to process each
input. Moving data between memory and the processor also contributes to
execution time and energy consumption. A reduction in the number of
operations does not therefore guarantee a proportional reduction in
execution time: the result also depends on the hardware and its
implementation~\cite{Sze2017EfficientDNN}.

In this thesis, the term \emph{lightweight} refers to models with lower
storage or computational requirements than a reference model. We will compare
the number of trainable parameters, the memory needed to store the model and
the estimated number of operations.
These quantities are considered together with detection performance
to assess which changes provide a useful reduction in resource requirements.

The reference model considered here is based on transformers. Their attention
mechanism combines information from different positions in an input
sequence~\cite{Vaswani2017Attention}, and this allows a model to use relationships
between observations at different times. TranAD applies this architecture
to anomaly detection in multivariate time series~\cite{Tuli2022TranAD}, while
TranAD+ builds on top of TranAD and evaluates a modified implementation on
several anomaly-detection benchmarks~\cite{Yermakov2025TranADPlus}.

We use TranAD+ as the reference model for our experiments. In the
implementation used here, the model receives a window of observations
and reconstructs its last sample, so that the reconstruction error is then used
to calculate an anomaly score. Effective detection requires this score
to distinguish nominal from anomalous observations. A lower reconstruction
error alone does not guarantee better anomaly detection, because errors
can decrease for both types of observation.

We study several ways to simplify the model: reducing the window size
limits the amount of past data processed for each observation; reducing
the feed-forward dimension makes the internal layers smaller;
architectural changes remove or reuse parts of the reference model;
quantization reduces the numerical precision
used to represent weights and to perform
calculations~\cite{Sze2017EfficientDNN}. For each change, we evaluate
both the reduction in resource requirements and its effect on
anomaly detection. We also compare TranAD+ with lightweight detectors
that do not use transformers and operate on simple statistics of the
signal.

\section{The TOAST project and thesis scope}
\label{sec:introduction-toast}

\ac{TOAST} is a research proposal at the University of Padua, involving
the Department of Physics and Astronomy and the Department of Information
Engineering. It aims to develop compact transformer models for onboard
anomaly detection in satellite telemetry. The project also plans their
implementation on embedded devices and the assessment of their robustness
against faults caused by radiation~\cite{TOASTExecutiveSummary}.

This thesis contributes to the model-development part of the proposal.
We first assess the reproducibility of TranAD+ on the benchmarks
considered in its original study~\cite{Yermakov2025TranADPlus},
and then investigate changes to the model and evaluate their effects
on detection performance and resource requirements. Finally, we consider
whether a detector without transformers can meet the same requirements,
since the choice of model for onboard use should follow from the evidence.

The main dataset is \ac{OPS-SAT}, which contains telemetry from an
experimental CubeSat operated by the \ac{ESA}. The data are organized
into single-channel signal fragments, each labeled as nominal or
anomalous~\cite{Ruszczak2025OPSSAT}. We use the measurements
within each fragment as a sequence of inputs with one feature.

The experiments in this thesis assess trained models through software
execution on the available computing platform. They provide information
about detection performance and computational requirements that can
support future onboard deployment. Implementation on embedded hardware
and radiation testing belong to later stages of the broader
\ac{TOAST} proposal~\cite{TOASTExecutiveSummary}.

\section{Objectives and research questions}
\label{sec:introduction-questions}

The objective of this thesis is to investigate how a transformer-based
anomaly detector can be simplified, and how these changes affect detection
performance and resource requirements. We start from TranAD+ as a reference
and study changes to its input window, internal dimensions, architecture
and numerical precision. The work addresses the following research questions:

\begin{enumerate}[label=\arabic*.]
  \item \textbf{Reference reproducibility.}
  How closely can we reproduce the published TranAD+ results using
  the available code, the same configurations and evaluation procedures?

  \item \textbf{Window size and feed-forward dimension.}
  How do changes to the input window and feed-forward dimension inside the transformer affect
  model performance?

  \item \textbf{Architectural simplification.}
  Can parts of the reference model be removed or shared while retaining
  detection performance?

  \item \textbf{Quantization.}
  How much can quantization reduce model storage, and how does it affect
  detection performance? What results are obtained with quantized weights
  and with integer matrix operations?

  \item \textbf{Evaluation protocol.}
  Do the metrics of the reference evaluation procedure distinguish
  models that separate anomalous from nominal observations from models
  that do not?

  \item \textbf{Lightweight alternatives.}
  Can a detector without transformers, based on simple statistics
  of the signal, provide better anomaly detection with lower resource
  requirements?
\end{enumerate}

The benchmark reproduction establishes a reference for the subsequent
comparisons. The evaluation procedure must also be considered, since
detection scores depend, for example, on how thresholds are selected.

These experiments test which resource reductions can be achieved
without losing detection capability, and whether they can also improve
detection performance. The direction and
size of each effect must therefore be established from the results.

For the model-comparison studies, repeated training runs are used
to assess sensitivity to random initialization and other random
operations during training. Comparisons use corresponding seeds,
and average performance is reported together with its variability:
this helps distinguish consistent effects from differences that
depend on an individual run.

\section{Contributions and organization of the thesis}
\label{sec:introduction-contributions}

As already stated, this thesis provides an experimental study of model simplification
for anomaly detection in satellite telemetry. We start from
the reproduction of TranAD+ benchmarks and its adaptation to
\ac{OPS-SAT}.

The subsequent experiments examine changes to window size,
feed-forward dimension, model architecture and numerical precision.
The architectural variants include a shared decoder, a single
encoder--decoder pass and an encoder followed by a reconstruction head.
Quantization is then evaluated using
the corresponding trained models as references. We also examine the
evaluation procedure itself and compare TranAD+ with lightweight
detectors without transformers.

The remaining chapters are organized as follows.
Chapter~\ref{chap:background} introduces time-series anomaly detection,
transformers, TranAD and TranAD+: it explains the model components
considered in the experiments and introduces lightweight detectors
without transformers.

Chapter~\ref{chap:datasets} presents the datasets, preprocessing,
training and evaluation procedures. It also describes anomaly scores,
threshold selection, point adjustment and the quantities used to assess
resource requirements.

Chapter~\ref{chap:replication} examines the reproduction of the
TranAD+ benchmarks and the adaptation to \ac{OPS-SAT}:
it discusses the differences between the reproduced and published
results and the checks performed to investigate them.

Chapter~\ref{chap:ablation} studies the effects of changing window size
and feed-forward dimension, both separately and together, while Chapter~\ref{chap:architecture}
presents the architectural variants
and Chapter~\ref{chap:quantization} shows quantization experiments, comparing detection performance
with resource requirements.

Chapter~\ref{chap:conclusions} discusses the main findings,
their limitations and possible developments within \ac{TOAST}.
